\chapter{Coset Graphs}

In this chapter we first show that for any group $G$, 
the transitive representations of $G$
together with $G$-homomorphisms between them
form a lattice isomorphic to a quotient of
the subgroup lattice of $G$.
%The primitive representations of $G$ are the co-atoms
%of this lattice.
Next we describe a construction of Sabidussi's
for vertex transitive graphs, and give a 
group theoretic characterization
of symmetric graphs.
Finally we consider quotients 
of symmetric graphs from a group theoretic perspective. 

%We describe three well-known methods of constructing ``extensions''
%of symmetric graphs -- the direct sum, the lexicographic product,
%and the Biggs covering graph construction.

\section{Transitive Permutation Groups}

% Coset Representation
\begin{definition}[Core of a subgroup]
If $H$ is a subgroup of $G$, then the \emph{core} of $H$ in $G$ is:
$$\Core_G(H) = \bigcap_{x \in G} x^{-1} H x $$
\end{definition}

\begin{proposition} If $H$ is a subgroup of $G$, 
then the core of $H$ in $G$ is a normal subgroup of $G$.
\begin{proof}
Since $gx$ runs over all the elements of $G$ as $x$ does, we have:
\begin{eqnarray*}
g^{-1}\Core_G(H)g & = & \bigcap_{x \in G} (gx) H (gx)^{-1} \\
                  & = & \bigcap_{x \in G} x H x^{-1} 
\end{eqnarray*}               
for any $g \in G$. 
\end{proof}
\end{proposition}

\begin{definition} [Coset Representation]
For $H$ a subgroup of $G$,
let $\Cos_G(H)$ denote the right cosets of $H$ in $G$.
We may define an action of $G$ on $\Cos_G(H)$ by:
$$ (Ha)^x = Hax $$
This is indeed an action, since:
\begin{eqnarray*}
{(Ha)^x}^y & = & (Hax)^y \\
           & = & (Haxy) \\
	   & = & (Ha)^{xy}
\end{eqnarray*}
for any ``$Ha$'' a coset of $H$ in $G$, and any $x,y \in G$.
We call a permutation group of the form $(G,\cos_G(H))$ a
{\it coset representation} of $G$.
\end{definition}

\begin{example}[Right Regular Representation]
For any group $G$, the \emph{right regular representation}
of $G$ is the permutation group $(G,G)$ with the action given by:
$$ g^{h} = gh. $$
This is a special case of a coset representation 
where $H$ is the trivial group $\{ 1 \}$.
\end{example}


\begin{proposition}
For any pair of groups $H \leq G$, the action of $G$
on the cosets of $H$ is transitive with kernel $\Core_G(H)$.
\begin{proof}
Let $Ha_1$ and $Ha_2$ be any two cosets of $H$ in $G$,
then:
$$(Ha_1)^{a_1^{-1}a_2} = Ha_2$$ 
so the action is transitive.

The stabilizer of the point ``$H$'' is $H$.
By Lemma 1, Chapter 1, 
the stabilizer of the point ``$Ha$'' is $a^{-1}Ha$.
If $g \in G$ is such that $g$ stabilizes every coset of $H$ in $G$,
then we must have $g \in a^{-1}Ha$ for each $a \in G$.
That is $g \in \Core_G(H)$.

\end{proof}
\end{proposition}


The next proposition tells us that every transitive representation
of a group $G$ is permutation equivalent to some coset representation.


\begin{proposition}
If $(G,\Omega)$ is a transitive permutation group
and $\alpha$ is some point of $\Omega$, then
$(G, \Omega)$ is permutation equivalent to $(G, \cos_G(G_\alpha))$

\begin{proof}
For each $\beta \in G$ Consider:
$$ S_\beta = \{ x \in G : \alpha^x = \beta  \} $$
If $\beta = \alpha^g$ then it's not hard to see
that $S_{\beta}$ is a right coset of $G_{\alpha}$:
$$ S_{\beta} = G_{\alpha} g $$
Define $\eta : \Omega \to \cos_G(G_a)$ by $\eta(\beta) = S_\beta$
and take $\varphi : G \to G$ to be the identity.
We have, for any $\beta \in \Omega$ and any $x \in G$,
if $\beta = \alpha^g$ then:
\begin{eqnarray*}
\eta(\beta^x) & = & S_{\beta^x} \\
              & = & S_{\alpha^{gx}} \\
	      & = & G_\alpha gx \\
	      & = & G_\alpha g^{\varphi(x)} \\
	      & = & \eta(\alpha^g)^{\varphi(x)} \\
	      & = & \eta(\beta)^{\varphi(x)}
\end{eqnarray*}
Thus $\eta$ and $\varphi$ establish the desired permutation equivalence.
\end{proof}
\end{proposition}
	
\section{Imprimitive Permutation Groups}
	
\begin{definition}
If $(G,\Omega)$ is a transitive permutation group and $\Delta$
is a subset of $\Omega$ then the {\it setwise stabilizer} of $\Delta$ is:
$$G_{\Delta} = \{ x \in G : \Delta^x = \Delta \}.$$
\end{definition}

\begin{lemma}
If $(G,\Omega)$ is a transitive permutation group and $\Delta$
is a block of imprimitivity, then
$G_{\alpha} \leq G_{\Delta} \leq G$.
for any $\alpha \in \Delta$.
\begin{proof}
Since $\alpha \in \Delta$, 
for any $g \in G$ we have $\alpha^g \in \Delta^g$.
Since $\Delta$ is a block of imprimitivity,
if $\alpha^g = \alpha$, then $\Delta^g \cap \Delta \neq \emptyset$
so $\Delta^g = \Delta$. That is $G_\alpha \leq G_\Delta$.
\end{proof}
\end{lemma}

\begin{proposition}
If $(G,\Omega)$ is a transitive permutation group and $\Delta$
is a block of imprimitivity, then $(G_{\Delta},\Delta)$
is a transitive permutation group.
\begin{proof}
Suppose that $\alpha, \beta \in \Delta$.
Since $(G,\Omega)$ is transitive, there is some $g \in G$
such that $\alpha^g = \beta$. 
Since $\alpha^g \in \Delta^g$ and $\beta \in \Delta$,
we must have $\Delta^g \cap \Delta \neq \emptyset$.
But $\Delta$ is a block of imprimitivity,
so this implies that $\Delta^g = \Delta$. 
That is $g \in G_{\Delta}$,
so $G_{\Delta}$ acts transitively on $\Delta$.
\end{proof}
\end{proposition}

\begin{proposition}
Suppose that $(G,\Omega)$ is a permutation group and $\alpha \in \Omega$.
For any subgroup $H$ such that $G_\alpha \leq H \leq G$, 
the set $\alpha^H = \{\alpha^h : h \in H \}$ is a block of imprimitivity.
\begin{proof}
Let $\Delta = \alpha^H$.
If $\Delta^g \cap \Delta \neq \emptyset $ for some $g \in G$
then there must be some $ x,y \in H $ such that
$\alpha^{xg} = \alpha^y$.
It follows that $ g \in x^{-1} G_\alpha y \leq H $
and $ \Delta^g = \Delta $.
That is $ \Delta^g \cap \Delta \neq \emptyset \Rightarrow \Delta^g = \Delta$,
so $\alpha^H$ is a block of imprimitivity. 
\end{proof}
\end{proposition}

\begin{lemma}
If $(G, \Omega)$ is a transitive permutation group
and $\alpha$ is some point of $\Omega$,
then $G_{\alpha^H} = H$
for any $H$ such that $G_\alpha \leq H \leq G$.
\begin{proof}
Let $\Delta = \alpha^H$.
Clearly $H \leq G_{\Delta}$.
To show that $H = G_\Delta$ it suffices to show that $G_\alpha$
and $G_\Delta$ have the same index in $H$.
Since $\alpha^H$ is an orbit of $H$ on $\Omega$,
the permutation group $(H,\alpha^H)$ is transitive.
Since $G_\alpha \leq H$, 
the stabilizer of the point $\alpha$ in $(H,\alpha^H)$
is $H \cap G_\alpha = G_\alpha$.
Thus by the orbit stabilizer theorem we have
$\Delta = [H : G_\alpha]$.

On the other hand, by proposition 11 and 12,
$\Delta$ is a block of imprimitivity
and $(G_{\Delta}, \Delta)$ is a transitive permutation group.
By lemma 3 we have $G_\alpha \leq G_\Delta$ 
so the stabilizer of the point $\alpha$ in $(G_\Delta, \Delta)$
is $G_\Delta \cap G_\alpha = G_\alpha$.
So again by the orbit stabilizer lemma we have
$|\Delta| = [G_{\Delta}: G_{\alpha}]$.
The result follows.

\end{proof}
\end{lemma}

\begin{lemma}
If $(G,\Omega)$ is a transitive permutation group,
$\Delta$ is a block of imprimitivity,
and $\alpha$ is any point of $\Delta$,
then $\alpha^{G_{\Delta}} = \Delta$.
\begin{proof}
Clearly $\alpha^{G_{\Delta}} \leq \Delta$.
Let $H = G_{\Delta}$.
We have $G_{\alpha} \leq H \leq G$
and $(H,\alpha^H)$ is a transitive permutation group
equivalent to $(H,\Cos_H(G_\alpha))$.
Since $(G_{\Delta}, \Delta)$ is also a transitive permutation group
equivalent to $(H, \Cos_H(G_\alpha))$
we must have $|Delta| = |alpha^H|$.
Since $\alpha^H \leq \Delta$, this implies that $\alpha^H = \Delta$.

\end{proof}

\end{lemma}

Let $(\cS, \leq)$ be the set of subgroups of $G$ containing $G_\alpha$,
partially ordered by the subgroup relation.
Let $(\cP, \subseteq)$ be the set of blocks of imprimitivity containing $\alpha$,
partially ordered by the subset relation.

\begin{proposition}
$(\cS, \leq)$ is order isomorphic to $(\cP, \subseteq)$

\begin{proof}
Let $\Phi : \cS \to \cP$ be given by
$\Phi(\Delta) = G_{\Delta}$
and let $\Psi : \cP \to \cS$ be given by
$\Psi(H) = \alpha^H$
Making use of lemma ?? and ??,
for any $H \in \cS$ we have:
$$ \Phi \circ \Psi(H) = \Phi(\alpha^H) = G_{\alpha^H} = H $$
and for any $\Delta \in \cP$ we have:
$$ \Psi \circ \Phi(\Delta) = \Psi (G_{\Delta}) = \alpha^{G_{\Delta}} = \Delta $$
so $\Phi$ and $\Psi$ are inverses.
To see that $\Phi$ is order preserving note that:
$$ G_{\Delta_1} \leq G_{\Delta_2} 
\Longleftrightarrow \alpha^{G_{\Delta_1}} \leq \alpha^{G_{\Delta_2}}
\Longleftrightarrow \Delta_1 \subseteq \Delta_2  $$
\end{proof}
\end{proposition}



\begin{proposition}
Let $(G,\Omega_1)$ and $(G,\Omega_2)$ be two transitive representations
of $G$. For any $\alpha \in \Omega_1$, $\beta \in \Omega_2$,
the permutation groups $(G,\Omega_1)$ and $(G,\Omega_2)$
are permutation equivalent if and only if there is an automorphism of $G$
carrying $G_\alpha$ to $G_\beta$.
\begin{proof}
Suppose that $(G, \Omega_1)$ and $(G, \Omega_2)$ are permutation equivalent.
Let $\varphi: G \to G$ and $\eta : \Omega_1 \to \Omega_2$ be the maps
establishing the equivalence. 
Let $\gamma = \eta(\alpha)$.
Clearly $G_\gamma = \varphi(G_\alpha)$.
By transitivity there must be some $g \in G$ such that $\beta = \gamma^g$.
Let $\vartheta_g : G \to G$ be the map: $x \mapsto g^{-1} x g$.
We have:
\begin{eqnarray*} 
G_\beta & = & g^{-1}G_\gamma g \\
        & = & \vartheta(G_\gamma) \\
	& = & \vartheta \circ \varphi (G_\alpha)
\end{eqnarray*}
So $\vartheta \circ \varphi$ carries $G_\alpha$ to $G_\beta$.

For the other direction, let $\psi$ be the map carrying $G_\alpha$
to $G_\beta$. By proposition 1,
$(G, \Omega_1)$ is permutation equivalent to $(G, \cos(G_\alpha))$
and $(G, \Omega_2)$ is permutation equivalent to $(G, \cos(G_\beta))$.
Take $\varphi = \psi$ and let 
$\eta: \cos(G_\alpha) \to \eta(G_\beta)$ be given by
$(G_\alpha) g \mapsto (G_\beta) \psi(g)$.
For any $x \in G$ we have:
\begin{eqnarray*}
\eta({(G_\alpha g)}^x) & = & \eta((G_\alpha) gx) \\
                     & = & (G_\beta) \psi(gx)  \\
		     & = & {(G_\beta) \psi(g)}^{\psi(x)} \\
		     & = & \eta((G_\alpha) g)^{\varphi(x)}
\end{eqnarray*}
thus $\varphi$ and $\eta$ establish a permutation equivalence between
$(G, \cos(G_\alpha)$ and $(G, \cos(G_\beta)$ and hence between
$(G, \Omega_1)$ and $(G, \Omega_2)$.
\end{proof}
\end{proposition}

Define an equivalence relation on the set of subgroups of $G$
as follows: $H_1 \sim H_2$ if and only if there is an automorphism of $G$
carrying $H_1$ to $H_2$. 
Let $\cZ$ denote the equivalence classes of this relation.
Define a partial order on $\cZ$
by $[H_1] \leq [H_2]$ if and only if there is some $\tilde{H_1} \in [H_1]$ and some
$\tilde{H_2} \in [H_2]$ such that $\tilde{H_1} \leq \tilde{H_2}$. 

The previous result tells us that, upto permutation equivalence,
the essential information about the transitive representations of $G$ 
is contained in $(\cZ, \leq)$. In particular all the transitive representations
of $G$ are to be found as quotients of the right regular representation.

\section{Sabidussi's Construction}

% Better make sure to discuss the difference between
% isomorphic as graphs and isomorphic as $G$-graphs.

We saw in the last section that
every permutation group is equivalent to
one of the form $(G,\cos_G(H))$.
Thus, upto isomorphism, transitive permutation groups are 
uniquely determined by pairs of groups $(G,H)$ 
with $H$ a subgroup of $G$.
In this section and the next
we shall see that upto isomorphism,
a symmetric graph is uniquely
determined by a triple $(G,H,a)$ 
where $H$ is a subgroup of $G$ 
and $a \in G \backslash H$ is an involution.

Given a triple $(G,H,a)$, the idea is to construct
a graph whose vertices are the cosets of $H$ in $G$.
We call such a graph a {\it coset graph}.
The idea of a coset graph originally goes back to Sabidussi
who was studying vertex transitive graphs.
Sabidussi's construction was essentially 
a generalization of the {\it Cayley graph} construction.

\begin{definition} (Cayley Graph)
Given a group $G$ and a subset $D \subseteq G$ the {\it Cayley graph},
$\Cay(G,D)$ is the directed graph with vertices the elements of $G$
and arc set $\{(x,y) : xy^{-1} \in D \}$
\end{definition}

A Cayley Graph will contain no loops,
provided that $1 \not\in D$.
A Cayley Graph is simple if and only if 
the set $D$ is closed under taking inverses,
and connected if and only if $D$ is a generating set for $G$.
The vertices are the elements of the group $G$,
and the action of $G$ on the vertices is permutation equivalent 
to the right regular representation of $G$.

Not every vertex transitive graph is a Cayley Graph.
Sabidussi's idea was to generalize Cayley's construction by considering the 
action of $G$ on the cosets of some non-trivial subgroup $H$.

\begin{definition} (Sabidussi Graph)
Given a group $G$, a subgroup $H$, and a set $D \subseteq G$, 
the Sabidussi graph, $\Sab(G,H,D)$, 
is the directed graph with vertex set $\cos_G(H)$ 
and arc set $\{(Hx,Hy) : xy^{-1} \in D\}$
\end{definition}

A Sabidussi graph will contain no loops
provided that $D \cap H = \emptyset$.
A Sabidussi Graph is simple if and only if
$D$ is closed under taking inverses
and connected if and only if $D \cup H$ generates $G$.
The vertices of the Sabidussi graph are  
the cosets of the subgroup $H$ in $G$,
and the action of $G$ on the vertices is $(G,cos(H))$.

\begin{proposition}[Sabidussi]
Sabidussi Graphs are vertex transitive.
Every vertex transitive graph is isomorphic to a Sabidussi Graph.
\begin{proof}
Let $\Gamma = \Sab(G,H,D)$ be any sabidussi graph
(simple and without loops).
The group $G$ has a natural transitive action on the vertices of $\Gamma$
given by $Hx^g = Hxg$.
To prove that $\Gamma$ is $G$-vertex transitive,
we must check that this action preserves the adjacency structure of $\Gamma$.

Suppose that $(Hx,Hy)$ is an arc of $\Gamma$,
so $xy^{-1} \in D$.
Since $xg(yg)^{-1} = xgg^{-1}y^{-1} = xy^{-1} \in D$
for any $g$, 
it follows immediately that $(Hx,Hy)^g = (Hxg,Hyg)$
is also an arc of $\Gamma$.
This proves the first part.

Now, let $\Gamma$ be any $G$-vertex transitive graph.
It follows that $(G, V \Gamma)$ is permutation
equivalent to $(G, \cos_G(H))$ for some $H < G$.
Let $\eta : \Omega \to \cos_G(H)$ be the map establishing
the permutation equivalence, and let $\mu = \eta^{-1}$

Let $N$ be the neighbourhood of the vertex $\mu(H)$, and let 
$$D = (\bigcup_{\alpha \in N} \eta(\alpha)) \diagdown H$$
We will show that $\Gamma \cong \Sab(G,H,\eta(D))$.

The map $\eta : V \Gamma \to \cos_G(H)$
establishes a permutation isomorphism between the
vertices of $\Gamma$ and the vertices of $\Sab(G,H,D)$.
Let $(\alpha, \beta)$ be any arc of $\Gamma$.
By $G$-vertex transitivity, we can find some $g \in G$
such that $\alpha^g = \mu(H)$ and $\beta^g \in E$.
Suppose that $\eta(\alpha) = Hx$ and $\eta(\beta) = Hy$.
We must check that $xy^{-1} \in D$

Since $\eta$ is a permutation isomorphism,
it follows that $\eta(\alpha^g) = Hxg$
and $\eta(\beta^g) = Hyg$.
Since $\beta^g \in E$ we must have 
$xg(yg)^{-1} = xgg^{-1}y = xy^{-1}\in D$.
The result follows.


\end{proof}

\end{proposition}

\section{A Group Theoretic Characterization of Symmetric Graphs}

Since symmetric graphs are vertex transitive,
every symmetric graph is a Sabidussi graph.
However, not every vertex transitive graph is symmetric,
so we expect there to be some extra conditions on
the subset $D$ in the symmetric case.

Recall from Chapter 1 that every symmetric graph 
is an orbital graph.
In the first section of this chapter, we saw that if $H < G$
is the stabilizer of some point in $\Omega$, then $(G,\Omega)$
is permutation equivalent to $(G,\cos_G(H))$.
We shall see in this section, the orbitals of $(G,\Omega)$
actually correspond to {\it double cosets} of $H$ in $G$

\begin{definition}
If $H$ is a subgroup of $G$, then a {\it double coset} of $H$ in $G$
is a subset of $G$ of the form
$HxH = \{ h_1 x h_2 : h_1, h_2 \in H \}$
for some $x \in G$.
\end{definition}

\begin{lemma} 
Each double coset of $H$ in $G$ is a union of right cosets of $H$.
The double cosets of $H$ in $G$ form a partition of $G$.
\begin{proof}
The first part is obvious since $HxH = \bigcup_{h \in H} Hxh$
for any $x \in G$.
For each element $g \in G$ it is clear that $g \in HgH$
so $\bigcup_{x \in G} HxH = G$.
If $HxH \bigcap HyH \neq \emptyset$ 
then $h_1 x h_2 = h_3 y h_4 $ for some $h_1,h_2,h_3,h_4 \in H$.
It follows that $y = h_1^{-1} h_3 x h_4 h_2^{-1}$ 
and $HyH = H h_1^{-1} h_3 x h_4 h_2^{-1} H = HxH$.
This shows that any two double cosets are equal or disjoint,
so the double cosets partition $G$ as claimed.
\end{proof}
\end{lemma}

\begin{proposition}[Lorimer]
For any permutationg group $(G,\Omega)$ with point stabilizer $H$, 
the orbitals are in natural bijection with the double cosets of $H$ in $G$.
The self-paired orbitals correspond to those 
double cosets which contain an involution.
\begin{proof}
From the first section in this chapter we know that $(G,\Omega)$
is permutation equivalent to $(G,\cos_G(H))$.
Recall from chapter 1 that the orbitals of a permutation group
are in bijection with the orbits of the point stabilizer.
We must show that orbit of $(H,\cos_G(H))$
corresponds to a double coset of $H$ in $G$.

Let $\kappa$ be the map $Hx \mapsto HxH$.
Suppose that $Hx$ and $Hy$ are in the same orbit
of $(H,\cos_G(H))$.
Then there is some $h \in H$ such that $Hxh = Hy$.
It follows that $y \in Hxh \subset HxH$.
That is $HyH = HxH$ and so $\kappa(Hx) = \kappa(Hy)$.

Conversly, suppose that $\kappa(Hx) = \kappa(Hy)$
for some pair of cosets $Hx$ and $Hy$.
Then $x \in HyH$ and so $x = h_1 y h_2$ for some 
$h_1, h_2 \in H$.
It follows that $Hx = Hyh_2$,
that is $Hx$ and $Hy$ lie in the same orbit of 
$(H,\cos_G(H))$.
Since $\kappa$ is surjective,
this proves the first part. 

Suppose that $a \in G$ is such that $a^2 = 1$ and $a \not\in H$.
Since $(H, Ha)^a = (Ha, Ha^2) = (Ha, H)$, it follows that 
the orbital $\Delta$ of $(G,\cos_G(H))$ 
containing $(H,Ha)$ is self-paired.
Clearly $a \in HaH = \kappa(Ha)$,
so the double coset associated with $\Delta$ contains an involution.

Conversely, if the double coset $HxH$ contains 
an involution $a$, then $HxH = HaH = \kappa(Ha)$.
Thus $HxH$ is associated with the orbital containing
$(H,Ha)$ which is self-paired.
This proves the second part.
\end{proof}
\end{proposition}

\begin{theorem} [Lorimer] For any pair of groups $H \leq G$ and any $a \in G$
such that $a \not\in H$ and $a^2 = 1$, the Sabidussi graph
$\Sab(G,H,HaH)$ is $G$-symmetric. Furthermore every symmetric graph
is of this form for some $G$, $H$ and $a$.
\begin{proof} 
Let $\Gamma = \Sab(G,H,HaH)$.
By proposition ?? $\Gamma$ is $G$-vertex transitive.
Since $HaH$ is an orbit of the action of $H$
on $\Cos_G(H)$, the graph $\Gamma$ is also $G$-locally transitive.

Suppose that $\Gamma$ is a $G$-symmetric graph.
Let $\alpha$ be any vertex, and let $H = G_\alpha$.
By proposition ??, the permutation group $(G, V \Gamma)$ 
is equivalent to $(G, \Cos_G(H))$.
Let $\eta : V \Gamma \to \Cos_G(H))$ be any map
inducing this equivalence.
Let $\beta$ be any neighbour of $\alpha$ in $\Gamma$
and let $a \in G$ be such that $(\alpha,\beta)^a = (\beta,\alpha)$.
Clearly $a$ is an involution and $\eta(\beta) = Ha$.

Let $\Gamma' = \Sab(G,H,HaH)$.
The map $\eta$ induces a permutation equivalence
between the vertices of $\Gamma$ and the vertices of $\Gamma'$.
To show that $\Gamma$ is isomorphic to $\Gamma'$
we must show that $\eta$ preserves adjacency.

Let $(\omega, \delta)$ be any arc of $\Gamma$.
Let $g \in G$ be such that $\omega^g = \alpha$.
Then $\delta^g$ is some neighbour of $\alpha$.
Thus there exists $h \in H$ such that $\delta^{gh} = \beta$.
Therefore $\eta(\delta^{gh}) = \eta(\beta) = Ha$.
Since $\eta$ is a permutation equivalence,
$\eta(\delta^{gh}) = \eta(\delta)^{gh}$,
so $\eta(\delta) = Ha^{{(gh)}^{-1}} = Hah^{-1}g^{-1}$.
Simmilarly, $\eta(\omega^g) = \eta(\omega)^g$
so $\eta(\omega) = H^{g^{-1}} = Hg^{-1}$.
Now $ah^{-1}g^{-1}g = ah^{-1} \in HaH$
so $(\eta(\omega),\eta(\delta))$ is an arc of $\Gamma'$.
The result follows.



\end{proof}
\end{theorem}

%\begin{proposition}
%$\Gamma = \Sab(G,H,HaH)$ is connected if and only if together
%$H$ and $a$ generalte $G$
%\end{proposition}

\begin{lemma}
$|HxH| = |H||H|/|x^{-1}Hx \bigcap H|$ Fix me
\begin{proof}
The mapping $HxH \to x^{-1}HxH$ given by $h_1 x h_2 \mapsto x^{-1} h_1 x h_2$
      is bijective, so $|HxH| = |x^{-1}HxH|$. 
      But for the rule for the product of two subgroups: 
      $|x^{-1}HxH| = |x^{-1}Hx||H|/|x^{-1}Hx \bigcap H|$,
      of course, $|x^{-1}Hx| = |H|$. and the result follows.
\end{proof}
\end{lemma}

\begin{proposition}
If $\Gamma$ is a $G$-symmetric graph isomorphic to the Sabidussi
graph $Sab(G,H,HaH)$ the the stabilizer of an arc of $\Gamma$
is isomorphic to $a^{-1}Ha \cap H$ and the valency of $\Gamma$
is $ |H|/|x^{-1}Hx \bigcap H| $.
\begin{proof}
Consider the arc $(H,Ha)$. The stabilizer of this arc is
the intersection between the stabilizer of the vertex ``$H$''
and the stabilizer of the vertex ``$Ha$''.
But the stabilizer of the vertex ``$H$'' is just $H$
and the stabilizer of the vertex ``$Ha$'' is $a^{-1}Ha$.
So the stabilizer of the arc is $a^{-1}Ha \cap H$.

Consider the action of $H$ on the neighbours 
of the vertex ``$H$''.
Since $\Gamma$ is symmetric, this action is transitive.
If ``$Ha$'' is any neighbour of ``$H$'',
then the action of $H$ on the neighbours of ``$H$''
is permutation equivalent to the action of $H$
on the cosets of the stabilizer of the arc $(H,Ha)$.
Since the stabilizer of the arc $(H,Ha)$ is $a^{-1}Ha \cap H$
it follows that the valency of $\Gamma$ is
$[H : a^{-1}Ha \cap H] = |H|/|x^{-1}Hx \bigcap H|$

\end{proof}

\end{proposition}

In some sense the valancy of a symmetric graph is
measuring the extent to which a subgroup fails to be normal.


\section{Quotient Graphs}

In the first section of this Chapter,
we determined, for any two transitive permutation groups
$(G, \Omega_1)$ and $(G,\Omega_2)$, 
the conditions under which there exists a
$G$-homomorphism from $(G,\Omega_1)$ to $(G,\Omega_2)$.

We saw that if $(G,\Omega_1)$
is permutation equivalent to $(G,\cos_G(H))$ 
for some $H < G$,
and $(G,\Omega_2)$ is permutation equivalent
to $(G,\cos_G(K))$ for some $K < G$,
then there is a $G$-homomorphism from $(G,\Omega_1)$
to $(G,\Omega_2)$ if and only if $H < K < G$.

In this section we give analogous conditions
under which there exists a $G$-homomorphism
from a $G$-symmetric graph $\Gamma$
to a $G$-symmetric graph $\Sigma$.

\begin{theorem} [Lorimer]
If $\Gamma$ is a $G$-symmetric graph isomorphic to $\Sab(G,H,HaH)$
and $\Sigma$ is a quotient of $\Gamma$,
then $\Sigma$ is isomorphic to $\Sab(G,K,KaK)$ for some $H < K < G$.
\begin{proof}
Let $\eta : V \Gamma \to \Cos_G(H)$ be the map inducing
the isomorphism between $\Gamma$ and $\Sab(G,H,HaH)$.
Let $\pi : V \Gamma \to V \Sigma$ be the map
inducing the homomorphism between $\Gamma$ and $\Sigma$.
Let $\alpha = \eta^{-1}(H)$,
let $B = \pi(\alpha)$ and let $\Delta = \pi^{-1}(B)$.

By proposition ?? $\Delta$ is a block of imprimitivity.
If $K = G_{\Delta}$ then the permutation group $(G, V \Sigma)$
is equivalent to $(G, \cos_G(K))$.
Let $\mu : V \Sigma \to \cos_G(K)$ be any map
inducing the equivalence.
If $\Sigma' = \Sab(G,K,KaK)$ then $\mu$ induces
a permutation equivalence between the vertices of $\Sigma$
and the vertices of $\Sigma'$.
We must check that it preserves adjacency. Blah blah blah...


\end{proof}
\end{theorem}

If $a \in K$ then the valency of the quotient is one.

%Butterflies in bloom. See how the grow from a small seed into such majestic beings?
